\subsection{Trajectory Generation}
\label{sec:generator-trajectories}

\subsubsection{Initial Histories and Random Innovations}
The initial history supplies the observations required by every lag leaf.
The AST represents the deterministic update; the executor adds the
innovation specified by the mechanism's noise law. These choices are retained
with the equation, and each generation family declares its own initialization
and innovation policies in Section~\ref{sec:generator-stability}.
The initial lag values and recursively generated observations have distinct
roles: the history is the starting state, and context/future windows select
their observations from the generated portion of the returned trajectory.

\subsubsection{Synchronous Multivariate Updates}
At target time $t$, a lag leaf reads $x_{j,t-\ell}$ with $\ell\geq1$.
Equivalently, from the common origin $\tau$, the first generated target uses
$t=\tau+1$. The shared executor
\path{src/tsfmxai/generators/stochastic.py} evaluates every target against the
same previous multivariate history, adds the supplied innovation, and then
appends all new values together. This synchronous update preserves the
intended lag relationships between channels. Model inputs and future targets
can subsequently be separated without changing the recorded equation.

\subsubsection{Multiple Trajectories from One Mechanism}
The equation structure, bound coefficients, and lag assignments stay fixed
while repeated trajectories receive separate initialization and innovation
streams. This distinguishes variation within one mechanism from variation
between sampled equations. The mechanism identifier connects those
trajectories to their shared equation and generation policy. Its executed
AST is also the source for the explanation references developed in
Section~\ref{sec:ground-truth-generation}.
